% META: DONE.1

\section{Форматы представления трёхмерных данных}

Для решения задачи восстановления истории построения необходимо работать с трёхмерными данными в различных форматах. В контексте данной работы ключевыми являются два формата: STEP (точное B-Rep представление) и STL (триангулированная сетка). Понимание их структуры, преимуществ и ограничений является необходимым условием для разработки эффективных методов анализа.

\subsection{Формат STEP: точное геометрическое представление}

STEP (Standard for the Exchange of Product Data) - это международный стандарт ISO 10303 для обмена данными о продукте между различными CAD-системами. Впервые опубликованный в 1994 году, STEP стал эталонным форматом для обмена 3D-данными в промышленности \cite{cadinterop2023}.

Основополагающей характеристикой формата STEP является использование B-Rep (Boundary Representation) - точного математического описания геометрии. В отличие от аппроксимирующих форматов, STEP хранит поверхности в аналитическом виде: NURBS-поверхности, аналитические кривые (окружности, линии), полную топологию (грани, рёбра, вершины) \cite{techsoft3d_step}; в частности, геометрия передаётся без потери точности, свойственной триангуляции \cite{bigrep2026}.

Структура STEP-файла включает несколько ключевых компонентов:
\begin{itemize}
\item Геометрическая информация: NURBS-поверхности, аналитические кривые, точные координаты.
\item Топологическая информация: описание того, как поверхности соединяются в единое тело (грани, рёбра, вершины).
\item Сборочная структура: иерархия компонентов в сборке.
\item Метаданные: информация о продукте, материале, допусках.
\end{itemize}

Современной версией стандарта является AP242 Edition 2 (2020 г.), которая добавляет поддержку семантической PMI (Product Manufacturing Information), GD\&T (Geometric Dimensioning and Tolerancing) и других атрибутов, необходимых для безбумажного производства \cite{rapidpipeline2026, 3dexperience_step}. STEP AP242 Edition 2 окончательно вытеснил устаревший формат IGES во всех сценариях, требующих семантической PMI и GD\&T \cite{3dviewstation2022}.

Для задачи восстановления истории построения STEP обладает как преимуществами, так и ограничениями. С одной стороны, точное B-Rep представление позволяет выполнять высокоточные геометрические вычисления, что критически важно для детерминированного определения параметров операций, а с другой стороны, STEP не хранит историю построения - это «замороженное» представление итоговой геометрии, что и создаёт проблему, которую призван решать данный проект.

\subsection{Формат STL: триангулированная сетка}

STL (Stereolithography) - это формат, разработанный для 3D-печати, но получивший гораздо более широкое распространение. STL-модель представляет собой пространственную триангуляционную сетку - поверхность, составленную из множества треугольников. Каждый треугольник (фасет) описывается координатами трёх вершин и (опционально) вектором нормали, и формат не поддерживает никакой дополнительной семантической информации - ни цветов, ни материалов, ни принадлежности треугольников к разным граням. Формат STL существует в двух вариантах: ASCII (текстовый, человекочитаемый) и бинарный (более компактный и распространённый) \cite{cvut2026}.

Главное преимущество STL - его простота и универсальность. STL поддерживается практически любым программным обеспечением, работающим с 3D-геометрией, не только в мире 3D-печати \cite{cccp3d2013}. Это делает его удобным базовым форматом для обмена геометрией между разными системами, практически гарантируя совместимость если в рамках подзадачи удаётся свести модель к этому формату. Однако у STL есть и существенные недостатки с точки зрения анализа геометрии. Во-первых, это аппроксимация - гладкая поверхность заменяется набором плоских треугольников, что приводит к потере точности. Во-вторых, отсутствие семантики - STL не различает грани, рёбра и вершины как топологические объекты; треугольники не сгруппированы в поверхности. В-третьих, зависимость от дискретизации - качество представления сильно зависит от параметра deflection (допуска триангуляции) \cite{flowvision2023}.

Для задачи машинного обучения на 3D-геометрии STL является одновременно и удобным, и проблематичным форматом. Удобным - потому что он предоставляет унифицированное представление в виде сетки, с которым могут работать стандартные нейросетевые архитектуры. Проблематичным - потому что триангуляция вносит шум, а отсутствие топологической информации затрудняет распознавание инженерной семантики (например, различия между скруглением и фаской).

\subsection{Формат OBJ: расширенное сеточное представление}

Помимо STL, в трёхмерном моделировании широко используется формат OBJ (Wavefront OBJ). Данный формат изначально был разработан для нужд компьютерной графики и анимации, однако благодаря своей гибкости и простоте нашёл применение и в инженерных задачах. В отличие от STL, формат OBJ позволяет описывать не только треугольники, но и многоугольные грани (четырёхугольники, пятиугольники и т.д.), что делает его более подходящим для представления тел-многогранников.

Структура OBJ-файла организована следующим образом: сначала последовательно перечисляются вершины как тройки пространственных координат (строки, начинающиеся с символа `v`), затем определяются грани (строки с `f`), в которых перечисляются индексы вершин (три и более), составляющих грань. Формат поддерживает также хранение нормалей, текстурных координат и групп объектов.

Важным преимуществом OBJ перед STL является возможность однозначного и без потери информации преобразования STL в OBJ. Более того, в формате OBJ легче обеспечить целостность водонепроницаемой поверхности, несмотря на возможные ошибки округления чисел с плавающей точкой, поскольку целостность проверяется исключительно по рёбрам и граням, и если два треугольника ссылаются на один и тот же индекс вершины, то гарантированно известно, что это одна и та же точка в пространстве. Это свойство упрощает задачи нормализации и анализа сетки.

Для нейросетевого анализа, в частности для архитектуры MeshCNN, OBJ является предпочтительным форматом, поскольку он сохраняет информацию о связности рёбер и граней, необходимую для вычисления инвариантных признаков (углы между гранями, двугранные углы). В рамках данной работы OBJ используется как основной формат для хранения размеченных сеток на этапе обучения и инференса.

\subsection{Сравнительный анализ форматов для задач машинного обучения}

Результаты анализа форматов представлены в сводной таблице \ref{tab:cmp-3d-format}

\begin{table}[h]
\centering
\begin{tabular}{|p{3.5cm}|p{4cm}|p{4cm}|p{3cm}|}
\hline
Характеристика & STEP (B-Rep) & STL (сетка) & OBJ (сетка) \\
\hline
Точность геометрии & Аналитическая, без потерь & Зависит от дискретизации & Зависит от дискретизации \\
\hline
Топологическая информация & Полная (грани, рёбра, вершины) & Отсутствует (набор независимых треугольников) & Частичная (явные индексы вершин, связность рёбер) \\
\hline
Семантическая информация & PMI, GD\&T, метаданные & Отсутствует & Отсутствует \\
\hline
Поддержка многоугольников & Да (любые грани) & Только треугольники & Да (треугольники и многоугольники) \\
\hline
Размер файла & Относительно компактный & Может быть большим при высокой детализации & Проме\-жу\-точный \\
\hline
Скорость обработки & Высокоточные вычисления требуют ресурсов & Быстрые визуализация и базовые операции & Умеренная \\
\hline
Пригодность для ML & Сложна для прямого использования (неоднородная структура) & Удобна, но нет топологии & Наиболее удобна для MeshCNN (сохраняет связность) \\
\hline
\end{tabular}
\caption{Сравнение форматов хранения 3D моделей}
\label{tab:cmp-3d-format}
\end{table}

Для гибридного подхода, предлагаемого в данной работе, оптимальным является комбинированное использование обоих форматов: STEP используется для точного детерминированного вычисления параметров и верификации гипотез на стороне CAD-ядра, а STL (или производные от него представления) - для нейросетевого анализа и распознавания геометрических паттернов. При этом для обучения и инференса нейросетевых моделей предпочтение отдаётся формату OBJ, как наиболее полно сохраняющему топологическую информацию о сетке.

